\documentclass{beamer}

\usepackage[utf8]{inputenc}
\usepackage[T5]{fontenc}
\usepackage[vietnamese]{babel}

\usepackage{amsmath,amssymb}
\usepackage{tikz}
\usepackage{listings}
\usepackage{array}
\usepackage{colortbl}
\usepackage{booktabs}
\usepackage{xcolor}
\usepackage{tcolorbox}

\usetheme{Madrid}
\usecolortheme{dolphin}

% --- màu ---
\definecolor{myblue}{RGB}{30,144,255}
\definecolor{mygreen}{RGB}{60,179,113}
\definecolor{myred}{RGB}{220,20,60}
\definecolor{mypurple}{RGB}{138,43,226}
\definecolor{myorange}{RGB}{255,140,0}
\definecolor{mycyan}{RGB}{0,206,209}
\definecolor{codebg}{RGB}{248,250,252}

\setbeamercolor{structure}{fg=myblue}
\setbeamercolor{frametitle}{fg=white,bg=myblue}

\tcbset{
colback=blue!5,
colframe=myblue,
arc=3mm
}

% macros
\newcommand{\hlblue}[1]{\textcolor{myblue}{#1}}
\newcommand{\hlgreen}[1]{\textcolor{mygreen}{#1}}
\newcommand{\hlred}[1]{\textcolor{myred}{#1}}
\newcommand{\hlpurple}[1]{\textcolor{mypurple}{#1}}
\newcommand{\hlorange}[1]{\textcolor{myorange}{#1}}


\title{Biểu Diễn Đồ Thị Trong Máy Tính}
\author{NHÓM 3}
\date{\today}

\begin{document}

\maketitle

\section{Khái niệm cơ bản}
Đồ thị là tập hợp các điểm biểu diễn mối quan hệ giữa các đại lượng. 
Ví dụ, hàm số:
\[
y = x^2
\]
sẽ tạo thành một đường cong parabol.

Máy tính biểu diễn đồ thị bằng cách:
\begin{itemize}
    \item Tính nhiều giá trị $(x, y)$
    \item Biểu diễn các điểm này trên màn hình
    \item Nối các điểm lại để tạo thành đường
\end{itemize}

\section{Cách máy tính biểu diễn đồ thị}

\subsection*{Bước 1: Chọn hệ trục tọa độ}
\begin{itemize}
    \item Trục hoành (Ox)
    \item Trục tung (Oy)
\end{itemize}

\subsection*{Bước 2: Lấy mẫu dữ liệu}
Chọn các giá trị của $x$ và tính:
\[
y = f(x)
\]

\subsection*{Bước 3: Vẽ điểm}
Mỗi cặp $(x, y)$ được biểu diễn bằng một điểm trên màn hình.

\subsection*{Bước 4: Nối các điểm}
Các điểm được nối lại để tạo thành đồ thị liên tục.

\section{Ví dụ minh họa}

Xét hàm số:
\[
y = x^2
\]

Bảng giá trị:

\[
\begin{array}{|c|c|}
\hline
x & y \\
\hline
-2 & 4 \\
-1 & 1 \\
0 & 0 \\
1 & 1 \\
2 & 4 \\
\hline
\end{array}
\]

Các điểm này được nối lại tạo thành một đường cong parabol.

\section{Các dạng đồ thị thường gặp}
\begin{itemize}
    \item Đồ thị hàm số (đường thẳng, parabol,...)
    \item Biểu đồ cột
    \item Biểu đồ đường
    \item Biểu đồ tròn
\end{itemize}

\section{Công cụ vẽ đồ thị}
\begin{itemize}
    \item Phần mềm: GeoGebra, Desmos
    \item Ngôn ngữ lập trình: Python, C++, Java
\end{itemize}

\section{Ứng dụng}
\begin{itemize}
    \item Học toán và trực quan hóa dữ liệu
    \item Phân tích dữ liệu
    \item Thiết kế đồ họa, game
    \item Khoa học và kỹ thuật
\end{itemize}

\end{document}